\section*{Sharp Boolean payment of the root signed-KL rows}

All logarithms in the GU functionals are base two.  Fix a Boolean
$f:\{-1,1\}^n\to\{-1,1\}$ and a physical coordinate $i$.  For
$A\subseteq[n]\setminus\{i\}$ put
\[
 a_A=\widehat f(A),\quad b_A=\widehat f(A\cup\{i\}),\quad
 m_A=a_A^2+b_A^2,
\]
\[
 C_A=\sum_{T\supseteq A}a_Tb_T,\quad
 M_A=\sum_{T\supseteq A}m_T,
\quad \alpha_A={2a_Ab_A\over m_A},\quad
\beta_A={2C_A\over M_A},
\]
with every zero-mass summand interpreted as zero.  The exact restriction
identity is
\[
 \Delta_iU=B_i-\Gamma_i+K_i,
 \qquad
 K_i=\sum_A m_A D_2\left({1+\alpha_A\over2}\middle\|
                           {1+\beta_A\over2}\right).       \tag{1}
\]

Write $K_i^{(0)}$ for the $A=\varnothing$ summand and $K_i^{(+)}$ for
the sum over $A\ne\varnothing$.

\paragraph{Theorem.}
For every Boolean $f$,
\[
 \boxed{\sum_i K_i^{(0)}
 \le {2\over\ln2}\sum_i\widehat f(\{i\})^2
 \le {2\over\ln2}I(f).}                              \tag{2}
\]
The coefficient $2/\ln2$ is optimal for this component.  Consequently,
write $c_0=2/\ln2$.  If some $a\ge c_0$ and $0\le b<1$ satisfy
\[
 \sum_iK_i^{(+)}
 \le (a-c_0)I+bU+\sum_i(\Gamma_i-B_i)                 \tag{3}
\]
on every tensor-prime Boolean function, then the universal GU inequality
$G\le aI+bU$ follows, and hence
\[
 H(f)\le\left(2+{2a\over1-b}\right)I(f)               \tag{4}
\]
for every Boolean $f$.

\paragraph{Boolean root normalization.}
The Fourier convolution identity for $f^2=1$ at the nonzero shift
$\{i\}$ is
\[
 0=\sum_{S\subseteq[n]}\widehat f(S)\widehat f(S\triangle\{i\})
  =2\sum_{A\subseteq[n]\setminus\{i\}}a_Ab_A.
\]
Therefore $C_\varnothing=0$ and $\beta_\varnothing=0$.  If
$a=\widehat f(\varnothing)$, $d=\widehat f(\{i\})$, and
$m=a^2+d^2$, then
\[
 K_i^{(0)}=m\left[1-h_2\left({1+2ad/m\over2}\right)\right]. \tag{5}
\]

\paragraph{Sharp scalar inequality.}
For $|x|\le1$, set
\[
 D(x)=1-h_2\left({1+x\over2}\right),\qquad
 R(x)=1-\sqrt{1-x^2}.
\]
Using natural logarithms,
\[
 (\ln2)D(x)={1\over2}\big((1+x)\ln(1+x)+(1-x)\ln(1-x)\big).
\]
Both sides vanish at zero.  For $0\le x<1$ their derivatives obey
\[
 {d\over dx}[(\ln2)D(x)]=\operatorname{artanh}x
 \le {x\over\sqrt{1-x^2}}=R'(x).                     \tag{6}
\]
Indeed the difference on the right minus the left vanishes at zero and
has derivative
\[
 {1\over(1-x^2)^{3/2}}-{1\over1-x^2}
 ={1-\sqrt{1-x^2}\over(1-x^2)^{3/2}}\ge0.
\]
Evenness and continuity at the endpoints prove
\[
 D(x)\le {1-\sqrt{1-x^2}\over\ln2}.                  \tag{7}
\]
Substitute $x=2ad/m$.  Since
\[
 \sqrt{1-x^2}={|a^2-d^2|\over a^2+d^2},
\]
(5)--(7) give the stronger pointwise estimate
\[
 K_i^{(0)}\le {2\min(a^2,d^2)\over\ln2}
 \le {2d^2\over\ln2}.                                \tag{8}
\]
Summing (8) proves (2).

For the $n$-bit point defect, $a=1-t$, $|d|=t$ with
$t=2^{1-n}$.  Dividing (5) by $d^2$ and sending $t\downarrow0$ gives
$2/\ln2$, so no smaller universal coefficient can replace $c_0$ in
(2).

\paragraph{GU implication.}
Summing (1), using (2), and then (3) yields
\[
 G\le c_0I+\sum_i(B_i-\Gamma_i)+\sum_iK_i^{(+)}
 \le aI+bU
\]
on tensor-prime functions.  The exact identities
$U(gh)=U(g)+U(h)$, $I(gh)=I(g)+I(h)$, and
$G(gh)=G(g)+G(h)$ for disjoint products extend the same pair to every
Boolean function.  Constants have $G=I=U=0$.  Finally, the restriction
induction implied by $G\le aI+bU$, together with the granularity theorem
$H\le2I+2U$, gives (4).

\paragraph{Scope.}
Only the empty upper-set rows have been paid.  The nonempty-row aggregate
in (3), including signed upper correlations, normalization, and possible
reuse across restrictions, remains open.  This theorem does not prove a
universal GU pair or full FEI.

